\chapter{The Ricci Flow of Special Geometries}

\section*{The Ricci flow}

The Ricci flow
\[
\frac{\partial}{\partial t} g = -2 \Rc
\]
with initial condition $g(0)=g_0$ and its cousin the normalized Ricci flow
\[
\frac{\partial}{\partial t} g
= -2 \Rc + \frac{2}{n}\frac{\int_{\M^n} R\, d\mu}{\int_{\M^n} d\mu}\, g
\]
are methods of evolving the metric of a Riemannian manifold $(\M^n,g_0)$.

\begin{remark}
\label{ch01-rem-manifolds-smooth-boundaryless}
\source{chow-knopf-ricci-flow:page-0015}
\dcref{Remark 1.1}
\group{topology-and-foundations}

In dimension $n=3$, the categories TOP, PL, and DIFF all coincide. In this volume, any manifold under consideration is assumed to be smooth. Unless explicitly stated otherwise, every manifold is assumed to be without boundary, so that it is closed if and only if it is compact.
\end{remark}

\section*{1. Geometrization of three-manifolds}

\begin{conjecture}
\label{ch01-conj-elliptization}
\source{chow-knopf-ricci-flow:page-0016}
\dcref{Conjecture 1.3 (Elliptization)}
\group{three-manifolds}

Let $\mathcal{N}^3$ be an irreducible orientable closed $3$-manifold of finite fundamental group $\pi_1(\mathcal{N}^3)$. Then $\mathcal{N}^3$ is diffeomorphic to a quotient $S^3/\Gamma$ of the $3$-sphere by a finite subgroup $\Gamma$ of $O(4)$. In particular, $\mathcal{N}^3$ admits a Riemannian metric of constant positive curvature.
\end{conjecture}

\begin{conjecture}
\label{ch01-conj-hyperbolization}
\source{chow-knopf-ricci-flow:page-0016}
\dcref{Conjecture 1.4 (Hyperbolization)}
\group{three-manifolds}

Let $\mathcal{N}^3$ be an irreducible orientable closed $3$-manifold of infinite fundamental group $\pi_1(\mathcal{N}^3)$ such that $\pi_1(\mathcal{N}^3)$ contains no subgroup isomorphic to $\mathbb{Z}\oplus\mathbb{Z}$. Then $\mathcal{N}^3$ admits a complete hyperbolic metric of finite volume.
\end{conjecture}

\begin{remark}
\label{ch01-rem-poincare}
\source{chow-knopf-ricci-flow:page-0016}
\dcref{Remark 1.5}
\group{three-manifolds}

A complete proof of the Elliptization Conjecture would imply the Poincar\'e Conjecture, which asserts that any homotopy $3$-sphere is actually a topological $3$-sphere.
\end{remark}

\section*{2. Model geometries}

\begin{proposition}
\label{ch01-prop-singer}
\source{chow-knopf-ricci-flow:page-0017}
\dcref{Proposition 1.7}
\group{model-geometries}

Any locally homogeneous metric on a simply-connected manifold is globally homogeneous.
\end{proposition}

\begin{proposition}
\label{ch01-prop-homogeneous-complete}
\source{chow-knopf-ricci-flow:page-0017}
\dcref{Proposition 1.8}
\group{model-geometries}

If $(\M^n,g)$ is homogeneous, then it is complete.
\end{proposition}

\begin{lemma}
\label{ch01-lem-model-geometry-to-homogeneous-space}
\source{chow-knopf-ricci-flow:page-0017,chow-knopf-ricci-flow:page-0018}
\dcref{Lemma 1.9}
\group{model-geometries}

Every model geometry $(\M^n,\G,\G_*)$ may be regarded as a complete homogeneous space $(\M^n,g)$.
\end{lemma}

\begin{proof}
\source{chow-knopf-ricci-flow:page-0017,chow-knopf-ricci-flow:page-0018}
\sketch Choose any $x\in \M^n$ and any scalar product on $T_x\M^n$. Average it over the compact isotropy group to obtain a $\G_x$-invariant scalar product, then transport that scalar product to every other tangent space using transitivity of the $\G$-action. The resulting metric is well defined, smooth, $\G$-invariant, and complete by Proposition~1.8.
\end{proof}

\begin{proposition}
\label{ch01-prop-isometry-group}
\source{chow-knopf-ricci-flow:page-0018}
\dcref{Proposition 1.10}
\group{model-geometries}

Let $(\M^n,g)$ be a smooth Riemannian manifold.
\begin{enumerate}
\item $\Isom(\M^n,g)$ is a Lie group and acts smoothly on $\M^n$.
\item For each $x \in \M^n$, the isotropy group $I_x(\M^n,g)$ is a compact subgroup of $\Isom(\M^n,g)$.
\item For each $x \in \M^n$, the linear isotropy representation $\lambda_x : I_x(\M^n,g) \to \Ogroup(T_x \M^n, g(x))$ is an injective group homomorphism onto a closed subgroup of the orthogonal group. In particular, $I_x(\M^n,g)$ is compact.
\end{enumerate}
\end{proposition}

\begin{corollary}
\label{ch01-cor-unique-isometry}
\source{chow-knopf-ricci-flow:page-0019}
\dcref{Corollary 1.11}
\group{model-geometries}

Given $x,y \in \M^n$ and a linear isometry $f : (T_x \M^n,g(x)) \to (T_y \M^n,g(y))$, there exists at most one isometry $\gamma \in \Isom(\M^n,g)$ with $\gamma(x)=y$ and $\gamma_* = f$. If $\M^n$ is compact, then $\Isom(\M^n,g)$ is compact.
\end{corollary}

\begin{lemma}
\label{ch01-lem-isotropy-isomorphic}
\source{chow-knopf-ricci-flow:page-0019}
\dcref{Lemma 1.12}
\group{model-geometries}

If $\Isom(\M^n,g)$ acts transitively on $\M^n$, then all isotropy groups $I_x(\M^n,g)$ are isomorphic.
\end{lemma}

\begin{proof}
\source{chow-knopf-ricci-flow:page-0019}
\sketch Given any $x,y \in \M^n$, choose an isometry $\gamma$ with $\gamma(x)=y$. Conjugation by $\gamma$ defines an isomorphism from $I_x(\M^n,g)$ to $I_y(\M^n,g)$, with inverse given by conjugation by $\gamma^{-1}$.
\end{proof}

\begin{remark}
\label{ch01-rem-model-geometry-ambiguity}
\source{chow-knopf-ricci-flow:page-0019}
\dcref{Remark 1.13}
\group{model-geometries}

The equivalence between model geometries and homogeneous models has an ambiguity: the same homogeneous model may arise from different transitive groups, and different isotropy sizes can lead to different compatible metrics. For the geometrization program, one restricts to maximal model geometries.
\end{remark}

\section*{3. Classifying three-dimensional maximal model geometries}

\begin{lemma}
\label{ch01-lem-unimodular-lie-group}
\source{chow-knopf-ricci-flow:page-0020}
\dcref{Lemma 1.14}
\group{maximal-model-geometries}

If $(\M^n,g)$ is a homogeneous model such that $\Isom(\M^n,g)$ is diffeomorphic to $\M^n$, then $\M^n$ is a Lie group. If moreover there exists a discrete subgroup $\Gamma$ of $\Isom(\M^n,g)$ such that $\M^n/\Gamma$ is compact, then $\M^n$ is unimodular.
\end{lemma}

\begin{proposition}
\label{ch01-prop-classification-maximal-models}
\source{chow-knopf-ricci-flow:page-0020,chow-knopf-ricci-flow:page-0021}
\dcref{Proposition 1.15}
\group{maximal-model-geometries}

Let $(\M^3,\G,\G_*)$ be a maximal model geometry represented by at least one compact $3$-manifold. Then exactly one of the following is true:
\begin{enumerate}
\item $\G_* \simeq \SO(3)$, in which case either
\begin{enumerate}
\item $\M^3$ is a Lie group isomorphic to $\SU(2)$;
\item $\M^3$ is a Lie group isomorphic to $\R^3$; or
\item $\M^3$ is diffeomorphic to hyperbolic $3$-space $\Hcal^3$.
\end{enumerate}
\item $\G_* \cong \SO(2)$, in which case either
\begin{enumerate}
\item $\M^3$ is a trivial bundle over a $2$-dimensional maximal model, so that
\begin{enumerate}
\item $\M^3$ is diffeomorphic to $S^2 \times \R$; or
\item $\M^3$ is diffeomorphic to $\Hcal^2 \times \R$;
\end{enumerate}
\item $\M^3$ constitutes a nontrivial bundle over a $2$-dimensional maximal model, so that
\begin{enumerate}
\item $\M^3$ is a Lie group isomorphic to $\nil$; or
\item $\M^3$ is a Lie group isomorphic to $\SL(2,\R)$.
\end{enumerate}
\end{enumerate}
\item $\G_*$ is trivial, in which case $\M^3$ is a Lie group isomorphic to $\sol$.
\end{enumerate}
\end{proposition}

\section*{4. Analyzing the Ricci flow of homogeneous geometries}

\begin{lemma}
\label{ch01-lem-curvature-offdiagonal}
\source{chow-knopf-ricci-flow:page-0023}
\dcref{Lemma 1.16}
\group{homogeneous-ricci-flow}

Suppose that $\{F_i\}$ is a Milnor frame for a left-invariant metric on a $3$-dimensional unimodular Lie group $G$. Then $\langle R(F_k,F_i)F_j,F_k\rangle = 0$ for all $k$ and any $i\neq j$.
\end{lemma}

\begin{proof}
\source{chow-knopf-ricci-flow:page-0023}
\sketch Using the Milnor-frame formulas for $\ad^*$ and the Levi-Civita connection, the mixed curvature component reduces to a product of orthogonal vectors, hence vanishes.
\end{proof}

\section*{5. The Ricci flow of a geometry with maximal isotropy $\SO(3)$}

\begin{proposition}
\label{ch01-prop-su2-collapse}
\source{chow-knopf-ricci-flow:page-0025,chow-knopf-ricci-flow:page-0026}
\dcref{Proposition 1.17}
\group{so3-isotropy}

For any $\epsilon \in (0,1]$ and any choice of initial data $\epsilon A_0, B_0, C_0 > 0$, the unique solution $g_\epsilon(t)$ to (1.5) exists for a maximal finite time interval $0 \leq t < T < \infty$. The metric $g_\epsilon(t)$ becomes asymptotically round as $t \nearrow T$.
\end{proposition}

\begin{proof}
\source{chow-knopf-ricci-flow:page-0025,chow-knopf-ricci-flow:page-0026}
\sketch After introducing $D=\epsilon A$ and ordering the variables so that $D\le C \le B$, one obtains differential inequalities forcing finite-time extinction and showing that $B,C,D$ stay comparable as $t$ approaches the singular time. This implies asymptotic roundness.
\end{proof}

\begin{remark}
\label{ch01-rem-round-sphere}
\source{chow-knopf-ricci-flow:page-0026}
\dcref{Remark 1.18}
\group{so3-isotropy}

When $B=C=D$ at $t=0$, this symmetry persists. It follows that $\SU(2)$ remains round and shrinks at the rate $D(t)=D_0-4t$.
\end{remark}

\begin{lemma}
\label{ch01-lem-common-time-interval}
\source{chow-knopf-ricci-flow:page-0027}
\dcref{Lemma 1.19}
\group{so3-isotropy}

There exists a time $T_0>0$ depending only on the initial data $(A_0,B_0,C_0)$ such that the solution $g_\varepsilon(t)$ of (1.5) exists for all $0\le t<T_0$ and all $0<\varepsilon\le 1$.
\end{lemma}

\begin{proof}
\source{chow-knopf-ricci-flow:page-0027}
\sketch The quantities $(AE)$ and $E/(BC)$ satisfy differential inequalities that yield a uniform lower bound for $BC$ on a time interval depending only on the initial data, and this in turn bounds the evolution of $A$ from below. Hence the solution persists on a common interval independent of $\varepsilon$.
\end{proof}

\begin{lemma}
\label{ch01-lem-collapse-discontinuity}
\source{chow-knopf-ricci-flow:page-0027,chow-knopf-ricci-flow:page-0028}
\dcref{Lemma 1.20}
\group{so3-isotropy}

If $(A_0,B_0,C_0)_\varepsilon$ is a family of initial data parameterized by $\varepsilon\in(0,1)$ such that $\lim_{\varepsilon\to 0}F(0,\varepsilon)>0$, then for all $t\in(0,T_0)$, $\lim_{\varepsilon\to 0}F(t,\varepsilon)=0$.
\end{lemma}

\begin{proof}
\source{chow-knopf-ricci-flow:page-0027,chow-knopf-ricci-flow:page-0028}
\sketch The evolution equation for $F$ has a negative term of size $Q/\varepsilon$ dominating the positive part uniformly on compact time intervals, so $F(t,\varepsilon)$ is forced to zero as $\varepsilon\searrow 0$.
\end{proof}

\section*{6. The Ricci flow of a geometry with isotropy $\SO(2)$}

\begin{proposition}
\label{ch01-prop-nil}
\source{chow-knopf-ricci-flow:page-0029,chow-knopf-ricci-flow:page-0030}
\dcref{Proposition 1.21}
\group{so2-isotropy}

For any choice of initial data $A_0,B_0,C_0 > 0$, the unique solution of (1.7) is given by (1.8). There exist constants $0 < c_1 \leq c_2 < \infty$ depending only on the initial data such that each sectional curvature $K$ is bounded for all $t \geq 0$ by $\frac{c_1}{t} \leq K \leq \frac{c_2}{t}$, and such that the diameter of any compact quotient $\mathcal{N}^3$ of $G^3$ is bounded for all $t \geq 0$ by $c_1 t^{1/6} \leq \diam \mathcal{N}^3 \leq c_2 t^{1/6}$. In particular, $\mathcal{N}^3$ is almost flat.
\end{proposition}

\begin{corollary}
\label{ch01-cor-nil-collapse}
\source{chow-knopf-ricci-flow:page-0030}
\dcref{Corollary 1.22}
\group{so2-isotropy}

On any compact nil-geometry manifold, the normalized Ricci flow undergoes collapse, exhibiting Gromov--Hausdorff convergence to $(\mathbb{R}^2,g_{\mathrm{can}})$, where $g_{\mathrm{can}}$ is the standard flat metric.
\end{corollary}

\section*{7. The Ricci flow of a geometry with trivial isotropy}

\begin{proposition}
\label{ch01-prop-sol}
\source{chow-knopf-ricci-flow:page-0031}
\dcref{Proposition 1.23}
\group{trivial-isotropy}

For any choice of initial data $A_0, B_0, C_0 > 0$, the unique solution of (1.10) exists for all positive time. For any $\epsilon > 0$, there exists $T_\epsilon \ge 0$ such that $\lvert A - \sqrt{A_0C_0} \rvert \le \epsilon$, $\lvert C - \sqrt{A_0C_0} \rvert \le \epsilon$, and $\left\lvert \frac{d}{dt} B - 16 \right\rvert \le \epsilon$ for all $t \ge T_\epsilon$. Moreover, there exist constants $0 < c_1 \le c_2 < \infty$ depending only on the initial data such that each sectional curvature $K$ is bounded for all $t \ge 0$ by $\frac{c_1}{t} \le K \le \frac{c_2}{t}$.
\end{proposition}

\begin{corollary}
\label{ch01-cor-sol-collapse}
\source{chow-knopf-ricci-flow:page-0031}
\dcref{Corollary 1.24}
\group{trivial-isotropy}

On any compact sol-geometry manifold, the normalized Ricci flow undergoes collapse, exhibiting Gromov-Hausdorff convergence to $\mathbb{R}$.
\end{corollary}

\begin{remark}
\label{ch01-rem-hyperbolization-progress}
\dcref{Remark 1.6}
\group{three-manifolds}
\source{chow-knopf-ricci-flow:page-0016}
Noteworthy progress toward the Hyperbolization Conjecture has come from topology.
\end{remark}

\begin{remark}
\label{ch01-rem-seifert-equivalence}
\dcref{Remark 1.2}
\group{three-manifolds}
\source{chow-knopf-ricci-flow:page-0016}
Seifert's original definition of a fiber space required the existence of a fiber-preserving diffeomorphism of a tubular neighborhood of each fiber to a neighborhood of a fiber in some quotient of $S^1 \times D^2$ by a cyclic group action. Epstein later showed that if a compact $3$-manifold is foliated by $S^1$ fibers, then each fiber must possess a tubular neighborhood with the prescribed property. Hence the simpler definition given in the text is equivalent to the original.
\end{remark}
